% 第7章 波形的发生和信号的转换
\chapter{波形的发生和信号的转换}

\section{正弦波振荡电路}
\subsection{概述}
正弦波振荡电路无需输入信号，通过正反馈产生一定频率、幅值的正弦波输出，自激振荡条件：
\[
\begin{cases}
\text{幅值条件：} |AF| = 1 \\
\text{相位条件：} \varphi_A + \varphi_F = 2n\pi \ (n=0,1,2,\dots)
\end{cases}
\]
核心组成：基本放大电路、正反馈网络、选频网络、稳幅环节。

\subsection{RC正弦波振荡电路}
选频网络为RC串并联网络（文氏桥），振荡频率：
\[ f_0 = \frac{1}{2\pi RC} \]
特点：适用于低频（1Hz~1MHz）振荡，输出波形失真小，频率调节方便。

\subsection{LC正弦波振荡电路}
选频网络为LC谐振回路，振荡频率：
\[ f_0 \approx \frac{1}{2\pi\sqrt{LC}} \]
\subsubsection{变压器反馈式LC振荡电路}
利用变压器绕组的同名端实现正反馈，适用于高频（1MHz以上）振荡，输出功率较大。

\section{电压比较器}
\subsection{概述}
电压比较器将输入电压与参考电压比较，输出高/低电平（非线性区工作，无负反馈），用于波形变换、越限报警等。

\subsection{单限比较器}
只有一个阈值电压 $U_T$，当 $u_I > U_T$ 时输出高电平，$u_I < U_T$ 时输出低电平（或相反），特点：结构简单，抗干扰能力差。

\subsection{滞回比较器（施密特触发器）}
有两个阈值电压（上限 $U_{T+}$、下限 $U_{T-}$），形成滞回特性，抗干扰能力强，阈值电压：
\[
U_{T+} = \frac{R_1}{R_1+R_2}U_{OH} + \frac{R_2}{R_1+R_2}U_{REF}, \quad U_{T-} = \frac{R_1}{R_1+R_2}U_{OL} + \frac{R_2}{R_1+R_2}U_{REF}
\]

\section{非正弦波发生电路}
非正弦波（方波、三角波、锯齿波）发生电路由滞回比较器和积分电路组成，核心是通过电容充放电改变阈值电压，实现自激振荡。

\section{例题参考}